\section{Managerial and Operational Implications}\label{sec:policy}

The evidence in Section~\ref{sec:experiments} is best read as tested simulation evidence
for DRT service design. It does not provide deployment rules or a production planning
tool. The implications below translate the model into bounded design considerations for
operators who are deciding whether and how to test bidirectional meeting-point service.

\subsection{Walking Tolerance as a Design Consideration}

Walking tolerance is the first constraint on bidirectional meeting-point consolidation.
When candidate pickup and dropoff points require too much walking, the binary-logit
response layer shifts otherwise feasible offers into choice rejection. Operators should
therefore treat walking-radius settings as local calibration parameters rather than
fixed rules. A practical pilot should observe pedestrian network quality, stop spacing,
weather exposure, safety conditions, and user willingness to walk before expanding the
service design.

\subsection{Fleet Deployment and Coverage Risk}

Fleet deployment affects whether meeting-point consolidation can translate into useful
service. Additional vehicles can reduce infeasibility and waiting burdens, but the
passenger-response layer can still constrain served share when walking or delay burdens
remain high. Operators should therefore evaluate vehicle-km per served trip together
with vehicle-km per original request, served share, choice rejection, and feasibility
rejection. A lower routing intensity signal is operationally meaningful only when the
coverage and rejection consequences are visible.

\subsection{Consolidation Trade-Off}

Bidirectional pickup and dropoff consolidation is not a universal replacement for
door-to-door service. It is a service-design option for conditions in which accepted
requests are spatially compatible and passenger burdens remain acceptable. Door-to-door,
single-sided pickup, and bidirectional service can coexist as operating modes. The
managerial question is which mode fits a demand density, fleet supply, walking tolerance,
and reliability target, not whether one mode is globally dominant.

\subsection{Passenger-Segment Monitoring}

The passenger-type results should be used as monitoring evidence for simulated
heterogeneity. They suggest that operators should inspect acceptance and burden patterns
by passenger segment when piloting meeting-point service. This monitoring can support
fallback design, opt-in rules, fare experiments, or differentiated service offers, but it
does not establish field equity effects. Passenger-segment monitoring should be
paired with empirical calibration before it is used for external reporting.

\subsection{Validation Limits and Diagnostic Evidence}

Several evidence layers support interpretation without becoming headline evidence.
Matched coverage is a diagnostic coverage-confounding control. Fixed accepted sets are a
diagnostic decomposition over fixed accepted sets. Gamma outputs are post-hoc welfare or
sensitivity accounting and do not alter routing, offer generation, or acceptance. The
Beijing-inspired synthetic grid is a tested simulation setting, not field validation. The
MILP comparison is a simplified ex-post diagnostic over fixed accepted sets, not a
complete dynamic routing standard. These validation limit statements should remain with
any operational interpretation of the findings.

\subsection{Calibration-Dependent VOT Interpretation}\label{subsec:vot-mapping}

The value-of-time interpretation of the utility coefficients is calibration-dependent.
The coefficients define the simulated passenger-type sensitivities used in this study;
they should not be interpreted as validated values for a particular population. Operators
who want to use the framework for local planning should estimate or stress-test walking,
waiting, in-vehicle time, and fare sensitivities from local stated-preference,
revealed-preference, or pilot acceptance data before drawing operational conclusions.
